\documentclass[12pt,a4paper,openany,oneside]{book}

% --- Paquetes ---
\usepackage[utf8]{inputenc}
\usepackage[spanish, es-noshorthands]{babel}
\usepackage{amsmath, amsfonts, amssymb}
\usepackage{graphicx}
\usepackage{geometry}
\usepackage{hyperref}
\usepackage{booktabs}
\usepackage{fancyhdr}
\usepackage{listings}
\usepackage{xcolor}
\usepackage{tikz}
\usepackage{pgfplots}
\usetikzlibrary{arrows.meta, positioning, shapes.geometric, calc}
\pgfplotsset{compat=1.18}

\geometry{margin=2.5cm}

% --- Colores y estilos para código Python ---
\definecolor{codegreen}{rgb}{0,0.6,0}
\definecolor{codegray}{rgb}{0.5,0.5,0.5}
\definecolor{codepurple}{rgb}{0.58,0,0.82}
\definecolor{backcolour}{rgb}{0.95,0.95,0.92}

\lstset{
    language=Python,
    backgroundcolor=\color{backcolour},   
    commentstyle=\color{codegreen},
    keywordstyle=\color{blue},
    numberstyle=\tiny\color{codegray},
    stringstyle=\color{codepurple},
    basicstyle=\ttfamily\small,
    breaklines=true,
    frame=single,
    showstringspaces=false
}

% --- Estilo de cabecera y pie ---
\pagestyle{fancy}
\fancyhf{}
\renewcommand{\headrulewidth}{0.4pt}
\renewcommand{\footrulewidth}{0.4pt}
\fancyhead[L]{\small Carlos César Sánchez Coronel}
\fancyhead[R]{\small Sesión 3: Operaciones con Tensores}
\fancyfoot[L]{\small Carlos César Sánchez Coronel}
\fancyfoot[C]{\thepage}

% --- Portada ---
\title{
    \textbf{Sesión 3} \\ 
    \Huge \textbf{OPERACIONES CON TENSORES} \\
    \Large Construyendo Algoritmos \\

}
\author{
    \small Por \\ \vspace{0.2cm}
    \Large \textsc{Carlos César Sánchez Coronel}
}
\date{2026}

\begin{document}

\maketitle
\tableofcontents
\addcontentsline{toc}{chapter}{Índice general}

% ------------------------------------------------------------------
% CAPÍTULO 3: SESIÓN 3 - OPERACIONES CON TENSORES I
% ------------------------------------------------------------------
\chapter{Operaciones con Tensores I: Construyendo Algoritmos}

En la sesión anterior aprendimos qué son los tensores y cómo manipular su forma. Ahora daremos el siguiente paso: las operaciones que nos permiten construir algoritmos de IA. Desde la simple suma de matrices hasta el producto punto que impulsa los mecanismos de atención en los grandes modelos de lenguaje. Estas operaciones son los ladrillos con los que se construyen las redes neuronales, las funciones de pérdida y los algoritmos de optimización.

\section{Operaciones elemento a elemento (element-wise)}

Las operaciones elemento a elemento se aplican de forma independiente a cada par de elementos correspondientes de dos tensores del \textbf{mismo shape}. Son la base de la computación paralela en GPUs.

\subsection{Suma y resta}
Dados dos tensores $\mathcal{A}, \mathcal{B} \in \mathbb{R}^{d_1 \times \dots \times d_k}$, su suma se define como:
\begin{equation*}
(\mathcal{A} + \mathcal{B})_{i_1,\dots,i_k} = A_{i_1,\dots,i_k} + B_{i_1,\dots,i_k}
\end{equation*}
De forma análoga para la resta.

\begin{lstlisting}[language=Python]
import numpy as np
A = np.array([[1,2],[3,4]])
B = np.array([[5,6],[7,8]])
print("Suma:\n", A + B)
print("Resta:\n", A - B)
\end{lstlisting}

\subsection{Producto Hadamard}
El producto Hadamard (o producto elemento a elemento) se denota a menudo con $\odot$:
\begin{equation*}
(\mathcal{A} \odot \mathcal{B})_{i_1,\dots,i_k} = A_{i_1,\dots,i_k} \cdot B_{i_1,\dots,i_k}
\end{equation*}
En NumPy se usa simplemente el operador \texttt{*}.

\begin{lstlisting}[language=Python]
print("Producto Hadamard:\n", A * B)
\end{lstlisting}

\begin{center}
\begin{tikzpicture}
\matrix[matrix of math nodes, left delimiter={[}, right delimiter={]}] (A) at (0,0) {
a_{11} & a_{12} \\
a_{21} & a_{22} \\
};
\matrix[matrix of math nodes, left delimiter={[}, right delimiter={]}] (B) at (3,0) {
b_{11} & b_{12} \\
b_{21} & b_{22} \\
};
\node[scale=2] at (1.5,0) {$\odot$};
\node[scale=2] at (4.5,0) {$=$};
\matrix[matrix of math nodes, left delimiter={[}, right delimiter={]}] (C) at (6,0) {
a_{11}b_{11} & a_{12}b_{12} \\
a_{21}b_{21} & a_{22}b_{22} \\
};
\end{tikzpicture}
\captionof{figure}{Producto Hadamard (elemento a elemento).}
\end{center}

\section{Broadcasting}

El \textbf{broadcasting} es una de las características más poderosas de NumPy (y de los frameworks de deep learning). Permite realizar operaciones entre tensores de diferentes formas, siempre que sean compatibles según ciertas reglas. La idea es que el tensor más pequeño se ``estira'' para que coincida con el más grande sin copiar datos explícitamente (lo que ahorra memoria y tiempo).

\subsection{Reglas de broadcasting}
\begin{enumerate}
    \item Se comparan las dimensiones de los dos tensores desde la última hacia la primera.
    \item Dos dimensiones son compatibles si:
        \begin{itemize}
            \item son iguales, o
            \item una de ellas es 1, o
            \item una de ellas no existe (se trata como 1).
        \end{itemize}
    \item Si se cumplen las reglas, el tensor con dimensión 1 se repite a lo largo de esa dimensión para igualar la del otro.
\end{enumerate}

\subsection{Ejemplos visuales}
\begin{itemize}
    \item \textbf{Escalar + vector}: Un escalar (shape vacío) se suma a cada elemento de un vector.
    \item \textbf{Vector + matriz}: Un vector fila (1, n) se suma a cada fila de una matriz (m, n); o un vector columna (m, 1) se suma a cada columna.
    \item \textbf{Imagen + escalar}: Sumar un valor a cada píxel (normalización).
\end{itemize}

\begin{center}
\begin{tikzpicture}
\draw[fill=blue!20] (0,0) rectangle (2,1) node[midway] {Matriz $3\times3$};
\draw[fill=red!20] (1,-1) rectangle (2,0) node[midway] {Vector $1\times3$};
\draw[->, thick] (1.5,-1) -- (1.5,0);
\node at (3,0.5) {Broadcasting: el vector se repite};
\draw[fill=red!20] (4,0) rectangle (5,1);
\draw[fill=red!20] (5,0) rectangle (6,1);
\draw[fill=red!20] (6,0) rectangle (7,1);
\end{tikzpicture}
\captionof{figure}{Broadcasting de un vector sobre las filas de una matriz.}
\end{center}

\subsection{Ejemplos en NumPy}
\begin{lstlisting}[language=Python]
import numpy as np

# Escalar + vector
a = 5
v = np.array([1,2,3])
print(a + v)  # [6 7 8]

# Vector + matriz (vector fila)
v = np.array([1,2,3])          # shape (3,)
M = np.array([[10,20,30],
              [40,50,60]])      # shape (2,3)
print(M + v)  # v se suma a cada fila

# Vector columna + matriz
v_col = np.array([[1],[2]])     # shape (2,1)
M = np.array([[10,20],
              [30,40]])          # shape (2,2)
print(M + v_col)  # v_col se suma a cada columna

# Ejemplo más complejo: tensor 3D + vector
T = np.random.rand(4,3,2)       # shape (4,3,2)
v = np.array([1,2])              # shape (2,) -> se puede broadcast a (4,3,2)
print((T + v).shape)              # (4,3,2)
\end{lstlisting}

\subsection{Aplicación en IA: Batch Normalization}
En \textbf{Batch Normalization}, se normaliza cada canal de un lote de imágenes. Se resta la media del canal y se divide por la desviación estándar. Estas operaciones usan broadcasting: la media tiene shape (1,1,c) y se resta de todo el lote (b,h,w,c).

\begin{lstlisting}[language=Python]
# Simulación simplificada
batch = np.random.rand(32, 32, 32, 3)  # (b, h, w, c)
media_por_canal = batch.mean(axis=(1,2), keepdims=True)  # (32,1,1,3)
batch_normalizado = batch - media_por_canal  # broadcasting
\end{lstlisting}

\section{Reducción (suma, media a lo largo de ejes)}

La reducción consiste en aplicar una operación (suma, media, máximo, etc.) a lo largo de uno o varios ejes, eliminando esas dimensiones.

\subsection{Suma total y suma por ejes}
\begin{equation*}
\text{sum}(\mathcal{X}) = \sum_{i_1,\dots,i_k} X_{i_1,\dots,i_k}
\end{equation*}
Para sumar a lo largo de un eje específico:
\begin{equation*}
(\text{sum}_j(\mathcal{X}))_{i_1,\dots,i_{j-1},i_{j+1},\dots,i_k} = \sum_{i_j} X_{i_1,\dots,i_k}
\end{equation*}

\subsection{Ejemplos en NumPy}
\begin{lstlisting}[language=Python]
M = np.array([[1,2,3],[4,5,6]])
print("Suma total:", M.sum())                # 21
print("Suma por filas (axis=1):", M.sum(axis=1))   # [6 15]
print("Suma por columnas (axis=0):", M.sum(axis=0)) # [5 7 9]
print("Media por columnas:", M.mean(axis=0))        # [2.5 3.5 4.5]
print("Máximo por filas:", M.max(axis=1))           # [3 6]
\end{lstlisting}

\begin{center}
\begin{tikzpicture}
\matrix[matrix of math nodes, left delimiter={[}, right delimiter={]}] (M) at (0,0) {
1 & 2 & 3 \\
4 & 5 & 6 \\
};
\node at (2,0.5) {$\xrightarrow{\text{sum(axis=0)}}$};
\matrix[matrix of math nodes, left delimiter={[}, right delimiter={]}] (S0) at (4,0) {
5 & 7 & 9 \\
};
\node at (6,0.5) {$\xrightarrow{\text{sum(axis=1)}}$};
\matrix[matrix of math nodes, left delimiter={[}, right delimiter={]}] (S1) at (8,0) {
6 \\ 15 \\
};
\end{tikzpicture}
\captionof{figure}{Suma a lo largo de diferentes ejes.}
\end{center}

\subsection{Aplicación en IA: Pérdida promedio en un batch}
Durante el entrenamiento, se calcula la pérdida para cada muestra en el batch y luego se promedia:
\begin{equation*}
\text{loss} = \frac{1}{B} \sum_{i=1}^{B} \ell(y_i, \hat{y}_i)
\end{equation*}
Esto es una reducción (media) sobre el eje del batch.

\section{Producto punto (dot product)}

El producto punto (o producto escalar) es una operación fundamental entre dos vectores que produce un escalar. Es la base de la similitud, de las proyecciones y de la atención.

\subsection{Definición}
Para dos vectores $\mathbf{u}, \mathbf{v} \in \mathbb{R}^n$:
\begin{equation*}
\mathbf{u} \cdot \mathbf{v} = \sum_{i=1}^{n} u_i v_i = \|\mathbf{u}\| \|\mathbf{v}\| \cos \theta
\end{equation*}
donde $\theta$ es el ángulo entre ellos.

\subsection{Producto matriz-vector y matriz-matriz}
El producto punto se generaliza a matrices. Para una matriz $\mathbf{A} \in \mathbb{R}^{m \times n}$ y un vector $\mathbf{x} \in \mathbb{R}^n$:
\begin{equation*}
(\mathbf{A} \mathbf{x})_i = \sum_{j=1}^{n} A_{ij} x_j
\end{equation*}
Para dos matrices $\mathbf{A} \in \mathbb{R}^{m \times n}$, $\mathbf{B} \in \mathbb{R}^{n \times p}$:
\begin{equation*}
(\mathbf{A} \mathbf{B})_{ij} = \sum_{k=1}^{n} A_{ik} B_{kj}
\end{equation*}

\begin{center}
\begin{tikzpicture}
\draw[fill=blue!20] (0,0) rectangle (2,2) node[midway] {$\mathbf{A}$};
\draw[fill=green!20] (3,0.5) rectangle (4,2.5) node[midway] {$\mathbf{x}$};
\node[scale=2] at (2.5,1.25) {$\cdot$};
\draw[fill=red!20] (5,0.5) rectangle (6,2.5) node[midway] {$\mathbf{y}$};
\node at (5.5,2.8) {Producto matriz-vector};
\end{tikzpicture}
\end{center}

\subsection{Propiedades}
\begin{itemize}
    \item \textbf{Conmutativo solo para vectores:} $\mathbf{u}\cdot\mathbf{v} = \mathbf{v}\cdot\mathbf{u}$.
    \item \textbf{Distributivo:} $\mathbf{A}(\mathbf{B}+\mathbf{C}) = \mathbf{A}\mathbf{B} + \mathbf{A}\mathbf{C}$.
    \item \textbf{Asociativo:} $(\mathbf{A}\mathbf{B})\mathbf{C} = \mathbf{A}(\mathbf{B}\mathbf{C})$ (cuando las dimensiones coinciden).
    \item \textbf{No conmutativo para matrices:} en general $\mathbf{A}\mathbf{B} \neq \mathbf{B}\mathbf{A}$.
\end{itemize}

\subsection{Similitud de coseno}
A partir del producto punto podemos definir la similitud de coseno:
\begin{equation*}
\text{sim}(\mathbf{u},\mathbf{v}) = \frac{\mathbf{u} \cdot \mathbf{v}}{\|\mathbf{u}\| \|\mathbf{v}\|}
\end{equation*}
Es una medida de orientación (el ángulo) independiente de la magnitud. Muy usada en embeddings y motores de búsqueda.

\subsection{Aplicación en Transformers: Atención}
El mecanismo de atención escalada por producto punto (dot-product attention) es el corazón de los transformers:
\begin{equation*}
\text{Attention}(Q,K,V) = \text{softmax}\left(\frac{QK^\top}{\sqrt{d_k}}\right) V
\end{equation*}
Aquí $Q$ y $K$ son matrices de consultas y claves. El producto $QK^\top$ calcula la similitud entre cada consulta y cada clave mediante productos punto. Esta operación es intensiva en cómputo y se beneficia de implementaciones eficientes (como las de las GPUs).

\begin{lstlisting}[language=Python]
# Simulación simplificada de atención
import numpy as np

def scaled_dot_product_attention(Q, K, V):
    d_k = Q.shape[-1]
    scores = np.matmul(Q, K.transpose(0,2,1)) / np.sqrt(d_k)
    weights = np.exp(scores) / np.sum(np.exp(scores), axis=-1, keepdims=True)
    return np.matmul(weights, V)

# Ejemplo: batch=2, seq_len=3, d_k=4
Q = np.random.rand(2,3,4)
K = np.random.rand(2,3,4)
V = np.random.rand(2,3,4)
output = scaled_dot_product_attention(Q, K, V)
print(output.shape)  # (2,3,4)
\end{lstlisting}

\section{Ejemplo integrado: De los conceptos a un algoritmo real}

Construyamos un pequeño algoritmo que use todas las operaciones vistas: normalización de un lote, cálculo de similitudes y reducción.

\subsection{Paso a paso}
\begin{enumerate}
    \item Tenemos un lote de embeddings de palabras (shape: batch, seq\_len, dim).
    \item Calculamos la media y desviación de cada embedding a lo largo de la dimensión de características (broadcasting).
    \item Normalizamos los embeddings (restamos media, dividimos por desviación).
    \item Calculamos la matriz de similitud de coseno entre todos los pares de palabras en cada secuencia (producto punto y normas).
    \item Para cada secuencia, obtenemos el valor máximo de similitud (reducción).
\end{enumerate}

\begin{lstlisting}[language=Python]
import numpy as np

# Datos simulados: batch de 2 frases, cada una de 3 palabras, embeddings de 5 dimensiones
batch = np.random.rand(2, 3, 5)
print("Shape batch:", batch.shape)

# 1. Normalización por características (broadcasting)
media = batch.mean(axis=2, keepdims=True)  # shape (2,3,1)
std = batch.std(axis=2, keepdims=True)     # shape (2,3,1)
batch_norm = (batch - media) / (std + 1e-8)

# 2. Para cada secuencia, calcular similitud de coseno entre cada par de palabras
def cos_sim_matrix(seq):
    # seq: (seq_len, dim)
    norms = np.linalg.norm(seq, axis=1, keepdims=True)  # (seq_len, 1)
    sim = np.dot(seq, seq.T) / (np.dot(norms, norms.T) + 1e-8)
    return sim

# Aplicar a cada secuencia del batch
sim_matrices = np.array([cos_sim_matrix(seq) for seq in batch_norm])
print("Shape de matrices de similitud:", sim_matrices.shape)  # (2,3,3)

# 3. Para cada secuencia, obtener la máxima similitud (excluyendo diagonal)
max_sim = []
for i in range(sim_matrices.shape[0]):
    # enmascarar diagonal (auto-similitud)
    sim = sim_matrices[i].copy()
    np.fill_diagonal(sim, -1)
    max_sim.append(sim.max())
print("Máxima similitud por secuencia:", max_sim)
\end{lstlisting}

\section{Conexión con futuros cursos}

\begin{itemize}
    \item \textbf{Deep Learning:} El producto punto es la operación fundamental en capas fully-connected ($\mathbf{W}\mathbf{x} + \mathbf{b}$). El broadcasting se usa en Batch Normalization y en la suma de sesgos.
    \item \textbf{NLP y Transformers:} La atención es producto punto entre consultas y claves. La similitud de coseno se usa para comparar embeddings.
    \item \textbf{Machine Learning:} En K-means, se calculan distancias (que involucran producto punto y normas). En SVM, el kernel lineal es un producto punto.
    \item \textbf{Visión por Computador:} Las convoluciones pueden verse como productos punto entre filtros y parches de la imagen.
    \item \textbf{Optimización:} El gradiente descendente requiere sumas y productos, y la reducción (promedio) sobre el batch.
\end{itemize}

\section{Resumen}

En esta sesión hemos aprendido las operaciones fundamentales que nos permiten construir algoritmos de IA:
\begin{itemize}
    \item Operaciones elemento a elemento (suma, resta, producto Hadamard).
    \item Broadcasting: cómo operar con tensores de diferentes formas.
    \item Reducción: sumar, promediar a lo largo de ejes.
    \item Producto punto: la base de la similitud y de las transformaciones lineales.
\end{itemize}
Cada una de estas operaciones tiene una interpretación matemática clara y una implementación eficiente en NumPy y en frameworks de deep learning. Dominarlas es esencial para entender y diseñar modelos de IA.

\end{document}